% from ega4-12-fr.tex
\oldpage[IV-3]{182}
  \item La fonction $y\mapsto\operatorname{prof}^{*}(\mathcal F_y)$ (10.8.1)
  est semi-continue inférieurement dans $Y$.
\end{enumerate}
\end{corollary}

\begin{proof}
(i) Il résulte de (12.2.1, (i)) appliqué à $\Phi$ égal au plus petit intervalle de
$\mathbf N$ contenant les dimensions des cycles premiers associés à
$\mathcal F_y$, que $z\mapsto\dim(\operatorname{Supp}(\mathcal F_z))$ est
semi-continue supérieurement au point $y$~; il résulte d'autre part de
(12.2.1, (ii)) appliqué à $\Phi$ égal à l'ensemble des dimensions des
composantes irréductibles de $\operatorname{Supp}(\mathcal F_y)$, que cette
même fonction est semi-continue inférieurement au point $y$, d'où la
conclusion. Les assertions (ii) et (iii) ne sont que d'autres formulations de
(12.2.1, (vi) et (xi)). Enfin, d'après (12.1.3), l'ensemble des $x\in X$ tels
que $\operatorname{prof}^{*}_{f(x)}(\mathcal F_{f(x)})\leq k$ est ouvert, et la
conclusion de (iv) en résulte par le même raisonnement qu'au début de (12.2.1).
\end{proof}

\begin{remarks}[12.2.3]
\label{IV.12.2.3-fr}
\begin{enumerate}[label=(\roman*)]
  \item On observera que pour (12.2.1, (ii)), on peut se dispenser de
  l'hypothèse que $f$ est \emph{propre}.

  \item Dans (12.2.1, (ii)), on ne peut remplacer la condition que l'ensemble
  $E(y)$ des dimensions des composantes irréductibles de
  $\operatorname{Supp}(\mathcal F_y)$ contienne $\Phi$, par la condition que
  $E(y)$ soit \emph{égal à} $\Phi$. En effet, soit $k$ un corps, et considérons
  l'espace projectif à 3 dimensions $P=\mathbf P^3_k=\operatorname{Proj}(C)$,
  où $C=k[t_0,t_1,t_2,t_3]$ ($t_i$ indéterminées)~; dans $P$, soit $X_0$ le
  sous-préschéma fermé «~réunion de la droite $X_1$ définie par $t_1=t_3=0$ et
  du plan $X_2$ défini par $t_2-t_3=0$~», ce qui décrit en termes géométriques
  le fait que $X_0$ est égal à $\operatorname{Proj}(C/\mathfrak a)$, où l'idéal
  gradué $\mathfrak a$ est égal à $\mathfrak b\mathfrak c$, avec
  $\mathfrak b=Ct_1+Ct_3$, $\mathfrak c=C(t_2-t_3)$~; considérons le morphisme
  $f_0:X_0\to X_1$ qui, en termes géométriques, est la «~projection de $X_0$
  sur $X_1$ à partir de la droite à l'infini $D$ définie par $t_0=t_2=0$~»~;
  sous forme algébrique, $f_0$ correspond à l'homomorphisme d'algèbres graduées
  \[
    k[t_0,t_2]\longrightarrow k[t_0,t_1,t_2,t_3]/\mathfrak a
  \]
  obtenu par passage au quotient à partir de l'injection canonique
  $k[t_0,t_2]\to k[t_0,t_1,t_2,t_3]$. Il est clair que $f_0$ est un morphisme
  projectif, donc \emph{propre}~; il en est donc de même de sa restriction
  $f:X\to Y$, où $Y=D_+(t_0)$ et $X=f_0^{-1}(Y)$~; pour voir que $f$ est
  \emph{plat}, il suffit de remarquer que $k[t_2]$ est un anneau principal et
  l'anneau
  \[
    C'=k[t_1,t_2,t_3]/(t_2-t_3)((t_1)+(t_3))
  \]
  un $k[t_2]$-module sans torsion ($\mathbf 0_{\mathrm I}$, 6.3.4). Pour tout
  $y\in Y$ distinct du point $y_0$ défini par $t_2=0$, $f^{-1}(y)$ a alors
  deux composantes irréductibles, de dimensions 0 et 1, tandis que
  $f^{-1}(y_0)$ n'a qu'une seule composante irréductible de dimension 1.

  \item Dans (12.2.1, (i)), on ne peut pas non plus remplacer la condition que
  l'ensemble $E'(y)\supset E(y)$ des dimensions des cycles premiers associés à
  $\mathcal F_y$ soit contenu dans $\Phi$ par la condition qu'il soit
  \emph{égal à} $\Phi$. Soient en effet $k$ un corps, $A$ l'anneau $k[t]$ de
  polynômes, $B$ le sous-anneau $k[t^4,t^3u,tu^3,u^4]$ de l'anneau de polynômes
  $k[t,u]$ ($t,u$ indéterminées)~; $B$ n'est pas intégralement clos, l'élément
  $t^2u^2$ appartenant à sa clôture intégrale mais n'étant pas dans $B$~; si
  $\mathfrak m$ est l'idéal maximal de $B$, engendré par $t^4,t^3u,tu^3$ et
  $u^4$, $\mathfrak m$ est idéal premier associé immergé de $A/At^4$. Posons
  $Y=\operatorname{Spec}(A)$, $X=\operatorname{Spec}(B)$, et soit $f:X\to Y$
  le morphisme correspondant à l'homomorphisme $A\to B$ qui transforme $t$ en
  $t^4$. Comme cet homomorphisme fait de $B$ un $A$-module sans torsion, $f$
  est plat ($\mathbf 0_{\mathrm I}$, 6.3.4). Si $y$ est le point de $Y$
  correspondant à l'idéal maximal $tA$, le préschéma $f^{-1}(y)$ est
  irréductible et de dimension 1, mais admet un cycle premier associé immergé
  de dimension 0~; au contraire, si $y'$ est le point générique de $Y$,
  $f^{-1}(y')$ n'a pas de cycle premier associé immergé, étant le spectre d'une
  $k(y')$-algèbre intègre de dimension 1. Dans cet exemple, $f$ est un morphisme
  affine non propre~; on peut en déduire aussitôt un exemple analogue où $f$
  est propre et plat en considérant $X$ comme un ouvert partout dense d'un
  schéma projectif sur $Y$ (\textbf{II}, 5.3.4 et 5.3.2), ou en procédant
  directement comme dans la Remarque (ii).

  Les Remarques (ii) et (iii) montrent la nécessité d'inclure dans (12.2.1,
  (iv)) la condition que $\mathcal F_y$ n'ait pas de cycle premier associé
  immergé.

  \item Dans (12.2.1, (xi)), on ne peut supprimer l'hypothèse que $f$ soit
  propre. En effet, soient $k$ un corps, $Y$ la «~droite affine~», spectre de
  l'anneau de polynômes $k[t]$ ($t$ indéterminée), $X$ le schéma \emph{somme}
  de $Y$ et de l'ouvert complémentaire du point fermé $y$ de $Y$ défini par
  $t=0$. Il est clair que le morphisme $f:X\to Y$ qui est égal aux injections
  canoniques dans chacune des composantes de $X$ est plat et que si l'on prend
  $\mathcal F=\mathcal O_X$, $\mathcal F_z$ n'a de cycle premier associé
  immergé pour aucun $z\in Y$. Toutefois, on voit aussitôt que l'on a $l(y)=1$
  et $l(z)=2$ pour tout $z\neq y$.

  \item Dans (12.2.1, (xi)), on ne peut pas non plus supprimer l'hypothèse que
  $\mathcal F_y$ soit sans cycle premier associé immergé. C'est ce que montre
  l'exemple de la Remarque (ii)~: on voit aussitôt en effet que l'on a, avec
  $\mathcal F=\mathcal O_X$, $l(y_0)=1$ et $l(y)=2$ pour $y\neq y_0$ dans $Y$.

  \item Nous ignorons si dans (12.2.1, (xi)) on peut remplacer l'hypothèse que
  $\mathcal F$ est $f$-plat par celle que $\operatorname{Supp}(\mathcal F)=X$
  et que $f$ est universellement ouvert. En reprenant la démonstration de
  (12.2.1.4), on voit (en utilisant aussi (14.3.6)) qu'il faudrait résoudre la
  question suivante~: Soient $A$ un anneau local noethérien, $M$ un $A$-module
  